% Mechanical relocation generated from the preservation ledger.
% Existing prose is included byte-for-byte from preserved blocks.

\section{Auxiliary Demand Baselines}
\subsection{Route-Level Models}

\subsubsection{Model Card: Route-Level IPF}

\paragraph{Formulation and required information.}
For one ordered route pattern, IPF finds a nonnegative upstream-to-downstream
vehicle-leg table matching the observed boarding and alighting marginals that
are explicitly marked as available.  It requires compatible marginals and a
strictly positive seed on the feasible triangular support.

\paragraph{Assumptions and limitations.}
The seed and entropy-like proportional adjustment determine the allocation not
fixed by the marginals.  With incomplete marginals the result is explicitly
underdetermined; even with complete marginals, many interior tables may share
the same margins.  Most importantly, a route boarding can be an initial
boarding or a transfer, so the fitted cells represent vehicle legs rather than
complete passenger journeys.  Route-level tables must not be stitched together
and interpreted as network OD demand without additional passenger-trajectory
information.

\paragraph{Appropriate use.}
This model is an inexpensive audit baseline, data-consistency check, and
possible initializer.  It is not a competing full-network journey-OD estimator.

For an ordered route with stops $1,\ldots,n$, let $x_{ij}\geq0$ for $i<j$ and
$x_{ij}=0$ otherwise. IPF alternately scales feasible rows and columns toward
the available boarding and alighting marginals,
\[
 \sum_{j>i}x_{ij}=b_i,
 \qquad
 \sum_{i<j}x_{ij}=a_j.
\]
Unobserved rows or columns are not constrained and remain dependent on the
seed. Complete marginals must be compatible with the triangular support; for
each prefix through $k$,
\[
 \sum_{j=1}^{k}a_j\leq\sum_{i=1}^{k-1}b_i.
\]
